% !TeX root=../../main.tex

\section{ساختار پایان‌نامه}

این پایان‌نامه در پنج فصل تدوین شده است که در ادامه، خلاصه‌ای از محتوای هر یک ارائه می‌شود:

\textbf{فصل اول (مقدمه و کلیات تحقیق):} به بیان مسئله، اهمیت کنترل نوسانات در جرثقیل‌های سقفی، مرور پیشینه پژوهش و تبیین اهداف و نوآوری‌های پروژه می‌پردازد.

\textbf{فصل دوم (مدل‌سازی دینامیکی و کنترل خطی):} شامل استخراج معادلات حرکت غیرخطی سیستم با استفاده از رهیافت اویلر-لاگرانژ، خطی‌سازی حول نقطه تعادل و طراحی کنترل‌کننده پایه \lr{LQR} است. همچنین، محدودیت‌های این رویکرد در مواجهه با شرایط اولیه بزرگ تحلیل می‌شود.

\textbf{فصل سوم (برنامه‌ریزی مسیر و کنترل پیش‌بین):} این فصل به توسعه برنامه‌ریز مسیر سینماتیکی و تدوین مسئله بهینه‌سازی برای کنترل‌کننده \lr{NMPC} اختصاص یافته است. همچنین نحوه بهره‌گیری از تغییر طول کابل جهت میرایی فعال نوسانات تشریح می‌شود.

\textbf{فصل چهارم (نتایج شبیه‌سازی و ارزیابی عملکرد):} عملکرد معماری کنترلی پیشنهادی را در محیط \lr{Simscape Multibody} مورد ارزیابی قرار می‌دهد. نتایج حاصل از سناریوهای مختلف حرکتی، تأثیر رویکرد غیرخطی را در مقایسه با روش‌های کلاسیک به اثبات می‌رساند.

\textbf{فصل پنجم (بحث و نتیجه‌گیری):} به جمع‌بندی دستاوردهای پژوهش، بررسی مزایا و محدودیت‌های روش پیشنهادی و ارائه پیشنهاداتی برای تحقیقات آتی می‌پردازد.